Pakistan's waste sector has no current national greenhouse-gas estimate built from Pakistani
activity data, and the international products that fill the gap disagree with each other by
amounts larger than the mitigation decisions they inform. We report the sector for
FY2023--24 at
\SI{37.45}{\mega\tonne} \COe{} (AR5, \SI{36.37}{\mega\tonne} on AR6) across three IPCC
categories: domestic wastewater (4D1) at \SI{24.46}{\mega\tonne}, solid-waste disposal (4A)
at \SI{9.17}{\mega\tonne} and open burning (4C2) at \SI{3.83}{\mega\tonne}. The current 10,000-draw conditional parameter ensemble gives \SIrange{25.4}{61.5}{\mega\tonne} CO$_2$e (2.5th--97.5th percentiles). It preserves the asymmetric effluent-N$_2$O factor treatment and the bounded waste-fate calculation; management alternatives and missing sources remain outside that interval. The same fate realisations support open-burning air pollutants, including \SI{63.85}{\kilo\tonne} primary PM$_{2.5}$, without double-counting generated material. The national annual core extends the disposal-history state and supported population inputs through FY2024--25; this note retains its detailed FY2023--24 analysis.

Two features distinguish the estimate from a conventional compilation. First, solid-waste
disposal and open burning are computed from one mass-balanced national waste-generation
figure rather than from separate literature: \SI{31.56}{\mega\tonne}
generated, \SI{50.3}{\percent} collected to disposal and the remainder uncollected, of which
\SI{60}{\percent} burns and the unburned \SI{6.27}{\mega\tonne} is carried as an explicit
unresolved fate with no emissions pathway applied.
Disposed plus uncollected equals generated, asserted with exact
equality rather than reconciled afterwards, so an argument about the collection rate moves
emissions between two categories without creating or destroying waste. Inventories that
estimate landfilling and open burning independently have no such constraint, and nothing then
forces their sum to be a quantity of waste the country actually produces.

Second, the largest category rests on an assumption the guidelines leave genuinely open for
Pakistani conditions. IPCC gives a methane correction factor of \num{0.5} for stagnant sewers
and \num{0.1} for discharge to a flowing river; Pakistani urban drainage has properties of
both. At \num{0.5} domestic wastewater is \SI{24.46}{\mega\tonne} and the sector sits at
\SI{141}{\percent} of EDGAR's year-matched Pakistani waste total; at \num{0.1} it is
\SI{15.39}{\mega\tonne} and the sector lands close to EDGAR. The physical description of the
drains supports the higher value and the only independent comparator supports the lower. We
report both, adopt \num{0.5} as central because the physical description is the stronger
evidence, and identify the measurement that would settle it: dissolved-methane sampling on
urban drains across a seasonal cycle. The physical description and the external inventory comparison play different roles; agreement with EDGAR is not used to select the central parameter. Seasonal measurements of drain methane could discriminate between the retained structural interpretations.
